\section{Candidate weights and their mass}\label{wt:section}

The ordinary bands $Q_j$ give an even number of prime slots and allow
the candidate to be split near a prescribed size. The smaller groups
$\mathcal P_g$ supply the graph marks and the odd parity restriction.
Separating the two families keeps these roles independent.

We now fix the prime groups used in the remainder of the proof. All
constants in this section depend only on the following fixed geometry
and on the function $\Psi$. In particular, they are independent of the
parameters $\kappa,b_1$ and of the proxy precision chosen later.

Choose a fixed even integer $r_0$ sufficiently large that
\[
 \frac{(r_0-1)c_1}{2}\ge c_2+1,
 \qquad c_2s'<\frac14,
\]
where
\begin{equation}\label{wt:geometry}
 c_1=1,\quad c_2=\frac65,\quad c_3=\frac32,\quad c_4=\frac95,
 \qquad s'=\frac1{r_0(c_1+c_2)},\qquad s_j=s'2^{-j}L.
\end{equation}
Let $j_x$ be the largest nonnegative integer with $c_1s_{j_x}\ge20T$,
and let $Q_j$ consist of the primes with
$c_1s_j\le\log p\le c_2s_j$, for $0\le j\le j_x$.
For each $j$ such that $[c_3s_j,c_4s_j]$ is contained in either
$[L^{1/10},L^{1/5}]$ or $[L^{3/10},L^{2/5}]$, take the primes whose
logarithms belong to $[c_3s_j,c_4s_j]$ as one group $\mathcal P_g$.
As before, $\mathcal P$ is the union of these groups and $n_*$ is the
part of $n$ supported outside $\mathcal P$.

These intervals are mutually disjoint, including between the $Q_j$
and the $\mathcal P_g$: indeed $c_2<c_3<c_4<2c_1$ and $c_4/2<c_1$.
The number of groups in a logarithmic range $[L^a,L^b]$ is
$(b-a)T/\log2+O(1)$. The prime number theorem and partial summation give,
uniformly in the indices,
\begin{equation}\label{wt:band-masses}
 V_{Q,j}:=\sum_{p\in Q_j}\frac1p=\log\frac65+o(1),
 \qquad
 V_g:=\sum_{p\in\mathcal P_g}\frac1p=\log\frac65+o(1).
\end{equation}
The uniformity follows because the smallest logarithmic endpoint tends
to infinity. Thus the group hypotheses of
Theorem~\ref{sh:correlation} hold with
$(a,b,c,d)=(1/10,1/5,3/10,2/5)$. Every prime in our construction is at
least $L^{20}$ once $x$ is sufficiently large. The classical prime
estimates used here and below are those in
\cite[Section ``Notation and preliminary estimates'']{SmoothShifted}.

Take $q_0=1/2$ and $\ell=1$ in the marked weights, so explicitly
\[
 W_1(n)=\prod_g \frac{q_0^{\omega_g(n)-1}\omega_g(n)}{V_g}.
\]
For a positive integer $v$, let $a(v)$ be the number of ordered prime
tuples with product $v$, with $r_0$ slots in every $Q_j$, divided by
$(r_0!)^{j_x+1}$. Repetitions are allowed. Fix a nonnegative smooth
function $\Psi$ supported in $[1,2]$ and positive on a nonempty open
subinterval of $(1,2)$. Define
\begin{equation}\label{wt:definition}
 \begin{gathered}
 \mathcal E(u)=(-1)^{\sum_{p\in\mathcal P}v_p(u)},\qquad
 A_0(u)=W_1(u)a(u_*),\\
 A(u)=A_0(u)\frac{1-\mathcal E(u)}2,
 \qquad X_A=\sum_u A(u)\Psi(u/x).
 \end{gathered}
\end{equation}
Here $v_p(u)$ denotes the exponent of $p$ in $u$.
If $e_p$ are the exponents of an integer in the support of $a$, then
\[
 a(v)=\prod_{j=0}^{j_x}\prod_{p\in Q_j}\frac1{e_p!},
 \qquad \sum_{p\in Q_j}e_p=r_0.
\]
In particular $0\le a(v)\le1$. The function
$k\mapsto kq_0^{k-1}$ is bounded on the nonnegative integers, and there
are $O(T)$ groups with $V_g$ bounded away from zero. Consequently
\begin{equation}\label{wt:pointwise}
 0\le A(u)\le A_0(u)\le L^{C_0}
\end{equation}
for a fixed $C_0$. The ordinary part of a supported integer has exactly
$r_0(j_x+1)$ prime factors counted with multiplicity. This number is
even, whereas the parity projection in $A$ selects an odd number of
group-prime factors. Thus
\begin{equation}\label{wt:parity}
 A(u)>0\quad\Longrightarrow\quad
 \Omega(u)\ \hbox{is odd},\qquad \Omega(2u)\ \hbox{is even}.
\end{equation}

\subsection{Harmonic measures and the group tail}

Set
\begin{equation}\label{wt:harmonic-masses}
 H_Q=\sum_v\frac{a(v)}v
     =\prod_{j=0}^{j_x}\frac{V_{Q,j}^{r_0}}{r_0!},
 \qquad
 H_P=\sum_{r:\,r_*=1}\frac{W_1(r)}r.
\end{equation}
Since $j_x+1=O(T)$, Equation~\eqref{wt:band-masses} gives
$L^{-C}\le H_Q\le L^C$ for a fixed $C$.
For one group introduce
\[
 Z_g(q)=\prod_{p\in\mathcal P_g}\left(1+\frac q{p-1}\right).
\]
The factor of $H_P$ associated with this group is
\begin{equation}\label{wt:group-mass}
 H_g=\frac{Z_g'(q_0)}{V_g}
 =\frac{Z_g(q_0)}{V_g}
       \sum_{p\in\mathcal P_g}\frac1{p-1+q_0}.
\end{equation}
This follows by expanding the Euler product: each occurring prime has
harmonic mass $\sum_{k\ge1}p^{-k}=1/(p-1)$, and differentiation chooses
one of the distinct prime divisors as its mark. Because $q_0\le1$, the
sum on the right of Equation~\eqref{wt:group-mass} is at least $V_g$.
The same sum is bounded above, and $Z_g(q_0)$ is bounded above, by
constants uniform in $g$. Hence
\begin{equation}\label{wt:HP-bounds}
 1\le Z_g(q_0)\le H_g\le C,
 \qquad
 H_P=\prod_gH_g,\qquad 1\le H_P\le L^C,
 \qquad
 \prod_{p\in\mathcal P}\left(1+\frac{q_0}{p-1}\right)\le H_P.
\end{equation}

Normalize $a(v)/v$ and $W_1(r)/r$ by $H_Q$ and $H_P$, respectively,
and draw the two parts independently. The $Q$-measure can equivalently
be sampled by drawing every ordered slot independently with law
$1/(pV_{Q,j})$ in band $j$. The group measure is a product of the
normalized measures belonging to Equation~\eqref{wt:group-mass}.
Write $V$ and $R$ for these two random integers and $U=VR$.

In one group the unnormalized harmonic mass of total multiplicity one
is exactly $V_g/V_g=1$. The contribution to multiplicity two from
two distinct primes is
\[
 \frac{2q_0}{V_g}\sum_{p<q\in\mathcal P_g}\frac1{pq}
 =\frac{q_0}{V_g}
     \left(V_g^2-\sum_{p\in\mathcal P_g}\frac1{p^2}\right).
\]
It is bounded below by a positive constant for sufficiently large $x$.
Since $H_g\le C$, each group therefore has both even and odd total
multiplicity with probability at least a fixed $\delta_0>0$.
Conditioning on all groups except one proves
\begin{equation}\label{wt:odd-mass}
 \mathbb P\bigl(\mathcal E(R)=-1\bigr)\ge\delta_0.
\end{equation}
The bound is independent of the number of groups.

\begin{lemma}[Tail of the group part]\label{wt:tail}
There is a fixed $C$ such that, for each fixed $\sigma>0$,
\[
 \sum_{\substack{r>x^\sigma\\r_*=1}}\frac{W_1(r)}r
 \le L^C\exp(-\sigma L^{3/5})
\]
for all sufficiently large $x$.
\end{lemma}

\begin{proof}
Put $\delta=L^{-2/5}$. The mass in one group after inserting $r^\delta$
is
\[
 \prod_{p\in\mathcal P_g}
      \left(1+\frac{q_0}{p^{1-\delta}-1}\right)
 \frac1{V_g}\sum_{p\in\mathcal P_g}
      \frac1{p^{1-\delta}-1+q_0}.
\]
Every group prime satisfies $p^\delta\le e$, so this expression is
bounded above by a fixed constant, using Equation~\eqref{wt:band-masses}.
Multiplying over $O(T)$ groups gives
$\sum_{r_* =1}W_1(r)r^{-1+\delta}\le L^C$.
For $r>x^\sigma$, the extra factor is at least
$x^{\sigma\delta}=\exp(\sigma L^{3/5})$, proving the assertion.
\end{proof}

\subsection{Localization at the target size}

\begin{lemma}[Total candidate mass]\label{wt:mass}
There is a fixed $C_A$ such that
\[
 X_A\asymp\frac{x}{L}H_QH_P,
 \qquad X_A\gg xL^{-C_A}.
\]
The implied constants and $C_A$ depend only on the fixed geometry and
on $\Psi$.
\end{lemma}

\begin{proof}
The exact probability identity is
\begin{equation}\label{wt:expectation}
 X_A=H_QH_P\,
 \mathbb E\left[U\Psi(U/x)\,
                   \1_{\{\mathcal E(R)=-1\}}\right].
\end{equation}
Leave one ordered slot in $Q_0$ free and condition on every other
variable. On the support of $\Psi(U/x)$, its logarithm must lie in an
interval of length $\log2$. Uniformly in the location of such an
interval, its normalized harmonic prime mass in $Q_0$ is $O(1/L)$.
For example, after intersecting with $Q_0$, the interval is contained
in $[y,2y]$ with $\log y\asymp L$, and the prime number theorem bounds
the reciprocal prime sum there by $O(1/L)$.
Since $U\le2x$ on the effective support, this proves the upper bound.

For the lower bound let
\[
 m_j=\frac{c_1+c_2}{2}s_j,
 \qquad \Delta=(c_2-c_1)s'L.
\]
The infinite midpoint sum satisfies
$r_0\sum_{j\ge0}m_j=L$, by the choice of $s'$.
Choose a fixed integer $J_0$ such that
\[
 \sum_{j>J_0}r_0c_2s_j<\Delta/24.
\]
For sufficiently large $x$ we have $j_x>J_0$. Restrict all slots in
$0\le j\le J_0$, except the free slot in $Q_0$, to fixed small
neighborhoods of their respective midpoints, choosing the widths so
that their total logarithmic deviation is less than $\Delta/24$.
There are only finitely many restrictions, each with harmonic
probability bounded below by a positive constant by the prime number
theorem. Their joint probability is therefore bounded below.

Independently, require $\mathcal E(R)=-1$ and
$\log R\le\Delta/12$. Equations~\eqref{wt:HP-bounds} and
\eqref{wt:odd-mass}, together with Lemma~\ref{wt:tail} applied to the
fixed positive number $\Delta/(12L)$, show that these two requirements
have probability at least $\delta_0/2$ for sufficiently large $x$.
No restrictions are needed on the remaining $Q$-slots. Their total
logarithm and the infinite midpoint sum for their bands both lie
between zero and $\Delta/24$. Consequently, if $B$ is the logarithm
of the product of all variables except the free prime, then on the
event just described
\[
 \left|B-(L-m_0)\right|<\Delta/6.
\]

Choose $1<\alpha<\beta<2$ with
$\Psi(t)\ge c_\Psi>0$ throughout $[\alpha,\beta]$.
For the free prime it is enough to require
\[
 \log p\in[L+\log\alpha-B,\ L+\log\beta-B].
\]
This interval has the fixed positive length $\log(\beta/\alpha)$.
Its endpoints lie strictly inside $[c_1s_0,c_2s_0]$ with distance
$\gg\Delta$ from the boundary, since $m_0$ is its midpoint and
$\Delta$ its full width. The prime number theorem and partial
summation give harmonic mass $\gg1/L$ for the interval, uniformly in
the conditioned variables. On it $U\ge\alpha x$ and
$\Psi(U/x)\ge c_\Psi$. Equation~\eqref{wt:expectation} now gives the
lower bound. Finally $H_Q\ge L^{-C}$ and $H_P\ge1$ imply the stated
fixed-power lower bound for $X_A$.
\end{proof}

\subsection{Squarefree predecessors}

\begin{lemma}[Squarefree loss]\label{wt:squarefree}
The mass of nonsquarefree predecessors satisfies
\[
 \sum_{\substack{u\ge1\\2u\ {\rm not\ squarefree}}}
       A(u)\Psi(u/x)=o(X_A/L).
\]
\end{lemma}

\begin{proof}
In one $Q$-band the probability that two given slots coincide is
\[
 \sum_{p\in Q_j}\frac1{p^2V_{Q,j}^2}
 \le\frac{L^{-20}}{V_{Q,j}}\ll L^{-20}.
\]
The number of pairs of slots within bands is $O(T)$, so the
probability of any repeated $Q$-prime is $O(TL^{-20})$.

For completeness, fix $p\in\mathcal P_g$ and write
$Z_{g,p}(q)=\prod_{q'\in\mathcal P_g,\ q'\ne p}
(1+q/(q'-1))$. The unnormalized mass of integers in this group that
are divisible by $p$ is
\[
 \frac{Z_{g,p}(q_0)+q_0Z_{g,p}'(q_0)}{V_g(p-1)}\ll\frac1p.
\]
The bound follows from the uniform group harmonic masses; division
by $H_g\ge1$ preserves it. Conditional on the set of distinct prime
divisors, the exponent of an occurring prime has distribution
proportional to $p^{-k}$, $k\ge1$. The conditional probability that
its exponent is at least two is therefore exactly $1/p$.
Thus the probability of a square from the group part is
\[
 \ll\sum_{p\in\mathcal P}\frac1{p^2}
 \le e^{-L^{1/10}}\sum_gV_g
 \ll T e^{-L^{1/10}}.
\]
All bands and groups are disjoint, and all their primes are odd.
It follows that $2U$ is not squarefree with probability
$O(TL^{-20})$ under the unrestricted harmonic measure.
Using Equation~\eqref{wt:expectation} and dropping parity and size
localization gives the upper bound
\[
 2x\|\Psi\|_\infty H_QH_P\,O(TL^{-20})
   \ll xH_QH_PT L^{-20}.
\]
By Lemma~\ref{wt:mass}, its ratio to $X_A/L$ is
$O(TL^{-18})=o(1)$, as required.
\end{proof}
